\scanpage{461}
\mspage{233}

Considérons comme élément de référence de $\unsure{\mathcal H^{-}|\mathcal G^{-}}$ l'élément
\[
(40)\qquad \tau=\tau_{\varnothing}\qquad
\text{(anti-involution \unsure{ss|s} pts fixes de $X_g$)}
\]
on trouve alors \inserted{par (32)} que pour $I'\subset I$
\[
(41)\qquad \tau_{I'}=\rho_{I'}\tau=
\left(\prod_{j\in I'}\rho_j\right)\cdot\tau\quad
\deleted{\illegible{1}}
\]
\margin{on voit \unsure{du reste|de suite}, comme $\tau^2=1$, $\tau_{I'}^2=1$, que
$\tau_{I'}=\tau_{I'}^{-1}=\tau\rho_{I'}^{-1}$
\unsure{donc} $\tau\rho_{I'}\tau=\rho_{I'}^{-1}$
\unsure{ce qui est $(42)$|ce qui $\Rightarrow(42)$}}
d'ailleurs on aura, pour \inserted{$\unsure{j|i}\in I$}
\[
\tau\rho_{\unsure{i|j}}\tau
=\tau_{\varnothing}\left(\tau_{\{\unsure{j|i}\}}\tau_{\unsure{\{\varnothing\}|\varnothing}}\right)
\tau_{\unsure{\{0\}|\{\varnothing\}}}
=\tau_{\varnothing}\tau_{\{j\}}=\rho_j^{-1}
\]
\[
(42)\qquad \tau\rho_{\unsure{i|j}}\tau^{-1}=\rho_{\unsure{i|j}}^{-1}.
\]

Ainsi $\mathcal G$ apparaît comme le produit semi-direct de $\mathcal H\simeq\mathbb Z^I$, et de
\deleted{\illegible{2}}
$\{\pm1\}\simeq\{1,\tau\}$ y opérant par la symétrie
$\rho\longmapsto\rho^{-1}$. Donc pour \uline{tout} $\rho\in\mathcal H$,
\unsure{\deleted{$\tau\rho$} on a}
$(\tau\rho)^2=1$, i.e. pour \unsure{tt|$\forall$}
$\sigma\in\unsure{\mathcal H^{-1}|\mathcal G^{-1}}$,
$\sigma$ \deleted{\illegible{1}} \uline{est une anti-involution} !
Comme $\sigma$ coïncide avec $\tau$ sur
$V_0=X_g-\bigcup_{\unsure{j|i}\in I}\operatorname{Int}(\unsure{\overline T_j|T_j})$,
on sait que l'ens des pts fixes de $\sigma$ est contenu dans
$\bigcup_{i\in I}\operatorname{Int}(T_i)$, et pour calculer l'indice de $\sigma$,
i.e. $\operatorname{card}\unsure{\pi_0|\Pi_0}(X_g^\sigma)$, il suffit de prendre la somme des
\margin{460}

\scanpage{462}

indices dans les $\unsure{\overline T_j|T_j}$. Or dans $\unsure{\overline T_j|T_j}$, on a \inserted{(29)}
$\rho_j(z,t)=(-z\exp\pi it,t)$, et par récurrence
\[
(43)\qquad \rho_j^{n_j}(z,t)=\left((-1)^{n_j}z\exp n_j\pi it,t\right)
\]
donc
\[
(44)\qquad \tau\rho_j^{n_j}(z,t)=
\left((-1)^{n_j+1}z\exp(n_j+1)\pi it,-t\right)
\]
et \deleted{\illegible{1}} $(z,t)$\inserted{$\in\unsure{\overline T_j|T_j}$} est pt fixe de
$\unsure{\tau|z}\rho_j^{n_j}$ \unsure{ssi|ss}
$t=0$, \deleted{\unsure{si|et}} et $n_j$ est \uline{impair}, donc
\[
(45)\qquad \unsure{\overline T_j|T_j}^{\tau\rho_j^{n_j}}=
\begin{cases}
\varnothing&\text{si $n_j$ pair}\\
S_j\times\{0\}&\text{si $n_j$ impair}
\end{cases}
\]

donc
\[
\begin{gathered}
(46)\qquad \text{Indice de }u=\tau\prod_j\rho_j^{n_j}\text{ est égal au cardinal}\\
\text{de l'ens. }I'_u\text{ des }j\in I\text{ tels que }n_j\\
\text{\unsure{sont|est} \unsure{impair|impairs}}
\end{gathered}
\]

Il est visible \inserted{par des raisons générales,} que $u$ est conjugué dans
$\unsure{\mathcal A_g|\widetilde{\mathcal A}_g}=\operatorname{Aut}(X_g)$
\inserted{(\unsure{par un|comme} élément de $\unsure{A_g^+|\widetilde A_g^+}$)}
à $\tau_{I'}=\rho_{I'}\tau=\tau\rho_{I'}^{-1}$ \unsure{ou aussi}
\deleted{
$\rho_{I'}^{-1}\tau=\tau\rho_{I'}$, \unsure{mais} \illegible{4}
\inserted{on a l'égalité \illegible{2} $\mathcal H$}
\unsure{dans le groupe} $\mathcal G$ (\unsure{trop petit}).
\illegible{4} $\rho\in\mathcal H$, \illegible{2}
\unsure{les questions}, d'une part
\[
\unsure{(\tau\rho)^{-1}}\cdot\rho_{I'}(\tau\rho)
=\unsure{\tau\rho\rho_{I'}}
\]
en introduisant des groupes de symétries supplémentaires, d'autre part en
\deleted{\unsure{l'agrandissant}, en} \unsure{complétant}
chaque $\unsure{\overline T_j|T_j}$ par une \unsure{configuration} plus
}
\margin{\deleted{
\[
\unsure{(\tau\rho)(\tau\rho')(\tau\rho)^{-1}}
=\underbrace{\tau\rho\tau}_{\unsure{\rho^{-1}}}
\underbrace{\rho'\rho^{-1}\tau}_{\unsure{\rho\rho'^{-1}}}
=\tau\rho^{\unsure{2}}\rho'^{-1}=\tau\rho''
\]
dès que $\rho^2=\rho'\rho''$, \unsure{c.à.d.|càd.}
\illegible{3} $\tau\rho'$ est \unsure{conjugué} \illegible{1}
$\unsure{\tau\rho''}$ dans $\mathcal G$, dès que
$\rho'\rho''\in\mathcal H^2$
}}
\margin{461}

\scanpage{463}
\mspage{234}

Mais si $\tau\rho',\tau\rho''\in\mathcal G^-$, $(\rho',\rho''\in\mathcal H)$, et si
$\rho\in\mathcal H$, la relation
\[
\rho(\tau\rho')\rho^{-1}=\tau\rho''
\]
équivaut à
\[
\underbrace{\tau\rho\tau}_{\rho^{-1}}\rho'\rho^{-1}=\rho''
\qquad\text{i.e.}\qquad \rho^2=\rho'\rho''^{-1}
\]
donc $\tau\rho'$ et $\tau\rho''$ sont conjugués par un élément de $\mathcal H$
\unsure{ssi|ss} $\rho'\rho''^{-1}\in\mathcal H^2$ -- ce qui précise l'observation précédente\ldots

La façon la plus riche en symétries simples pour disposer les $2g$ trous antipodiques
$\unsure{D'_j|D_j}$ me semble la suivante. On \deleted{\illegible{1}} considère
la sphère euclidienne, avec \unsure{l'}action du groupe dihédral
$\unsure{\mathbb D_{2g}|\mathbf D_{2g}}$ -- par exemple, quand c'est la sphère de Riemann qui est
considérée comme riemannienne, par le choix de l'antipodisme comme étant
\[
(47)\qquad \tau z=-\frac{1}{\overline z}
\]
on prend l'action type du groupe dihédral (avec comme pôles les points $0,\infty$, et
comme équateur le cercle unité $\mathbb U=\{z\in\mathbb C\mathbin{:}|z|=1\}$),
\deleted{par} en écrivant
$\unsure{\mathbb D_n|\mathbf D_n}$ comme $\subset\operatorname O(\deleted{\unsure{3}}2,\mathbb R)$,
comme un produit $\frac12$ direct de
\deleted{\unsure{$\mathbb Z/2\mathbb Z$}}\inserted{$\{\pm1\}$} par $\mu_n(\mathbb C)=\mu_n$,
le couple $(\zeta,\alpha)$ $(\zeta\in\mu_n,\alpha\in\deleted{\illegible{1}}\inserted{\{\pm1\}})$ opérant par
\[
(48)\qquad u_{\zeta,\alpha}(z)=\zeta z^\alpha
\]
\margin{462}

\scanpage{464}

On peut d'ailleurs l'élargir en un groupe
$\mathbb D_n^\sim=\mathbb D_n\times\mathbb Z_{/2}$,
où le deuxième facteur $\mathbb Z_{/2}$ est engendré par l'antipodisme (47)
(qui commute à $\mathbb D_n$ --- de $\overline{\mathrm m}$ d'ailleurs
que l'anti-involution $z\longmapsto\frac{1}{\bar z}$, qui a comme
ens des pts fixes $\mathbb U$ et n'est autre que la symétrie p.r. à
l'équateur, leur composé $z\longmapsto-z$ étant la symétrie p.r. à
l'axe des pôles \ldots) donc de regarder les transformations
\[
u_{\zeta,\alpha,\beta}\deleted{\illegible{1}}
=\unsure{\zeta}u_{\zeta,\alpha}\tau^\beta
\qquad
(\zeta\in\mu_n,\ \alpha\in\{\pm1\},\ \beta\in\mathbb Z_{/2})
\]
donc
\[
u_{\zeta,\alpha,\beta}z=
\begin{cases}
u_{\zeta,\alpha}(z)=\zeta z^\alpha&\mathrm{si}\ \beta=0\\
u_{\zeta,\alpha}(-\frac{1}{\bar z})=-\zeta\bar z^{-\alpha}
&\mathrm{si}\ \beta=1
\end{cases},
\]
mais on se \unsure{rappellera} que $\tau$ renverse l'orientation,
donc \deleted{\illegible{1}}\inserted{\unsure{est à|à}}
\deleted{\unsure{plus}}\inserted{\unsure{priori}} avec
réserve pour ce qui concerne le "transport de structure" dans
la situation présente.

Ici $n=2g$, $\mathbb D_{\unsure{2g|n}}$ est d'ordre $4g$, et
$\mathbb D_n^\sim$ d'ordre $8g$. On prend sur l'équateur une
trajectoire $\unsure{\mathcal J|\mathbb J}$, de $\mu_n=\mu_{2g}$, p.ex.
\unsure{les|ici} \unsure{justement}
\deleted{\unsure{les}} ($\mu_{2g}$
\inserted{\unsure{l'ens} des racines \unsure{$2g$ ièmes} de l'unité}
lui-$\overline{\mathrm m}$, \unsure{vu}\deleted{\illegible{1}}
\unsure{pour} sous-ens. de la sphère) comme l'ens des "centres"
des disques $D'_i$. On choisit un disque $D'_0$ autour du
\unsure{point|pt} $P_0$ (assez petit pour ce qui va suivre), et on
prend les \unsure{transformés|transformées} de $D'_0$ par les
$\zeta\in\mu_n$.

\margin{il faut \unsure{simplemt} que $D'_0$ ne
\unsure{rencontre} pas $\unsure{\zeta_0|\zeta}D'_0$,
\unsure{où} $\zeta=\exp\frac{2i\pi}{2g}$}

\margin{463}

\scanpage{465}
\mspage{235}

Ces choix étant faits, la surface $X_g$ est déterminée sans
\unsure{ambiguité|ambiguïté}, et le groupe $\mathbb D_{2g}$ y opère par transport de
structure, en permutant entre eux les $g$ cylindres $T_i$,
\unsure{correspond.ts|correspondants} aux éléments de
\inserted{$\unsure{\mathcal J|\mathbb J}/\tau=$}$\unsure{\mathcal J|\mathbb J}/\{\pm1\}$ i.e.
aux paires d'éléments antipodiques de $\unsure{S|\mathcal J}$,
i.e. aux "diagonales" du polygone à $2\unsure{g|\mu}$ côtés
qu'ils \deleted{\unsure{forment}}\inserted{déterminent}
sur l'équateur. Le groupe $\mathbb D_{2g}$ normalise
\deleted{\illegible{1}} le groupe $\mathcal G$, de façon précise
on aura, pour $u\in\mathbb D_{2g}$
\[
\tag{49}
u\tau_{I'}u^{-1}=\tau_{u(I')}
\]
pour $I'\subset I=\unsure{\mathcal J|\mathbb J}/\tau$, en tenant compte de
l'opération de $\mathbb D_{2g}$ sur $I$. On aura donc en
particulier, désignant par $\tau_g$\inserted{$=\tau_\varnothing$}
l'extension \margin{\deleted{(50)}} de l'antipodisme de la sphère $X_0$
(ou plutôt de $\unsure{\tau}\vert V_0$) en un antipodisme
sans pts fixes de $X_g$, noté $\tau$ précédemment:
\[
\tag{50}
u\tau_g u^{-1}=\tau_g
\qquad\text{i.e. $\tau_g$ commute à l'action de $\mathbb D_n$}
\]
\[
\deleted{\unsure{(51)}}\qquad
u\rho_j u^{-1}=\rho_{u\unsure{(}j)}
\]

On peut donc dire que $\mathbb D_{2g}$ opère sur $\mathcal G$,
d'où un produit $\frac12$-direct $\mathbb D_{2g}.\mathcal G$,
qui opère donc sur $X_g$.
Quant à la question de prolonger de $\overline{\mathrm m}$
l'action \inserted{sur $X_0$,} de $\mathbb D_{2g}^\sim$
\deleted{sur $X_0$} \unsure{tout entier} en une action sur $X_g$,

\margin{464}

\scanpage{466}

ça a été fait sans crier gare, à $\tau_{X_0}$ correspondant
naturellement $\tau_{X_g}=\tau_g$, qui a en effet le bon goût
de commuter à l'action de $\mathbb D_{2g}$.
[Il pourrait sembler plus naturel, il est vrai, dans un
esprit de "transport de structure (\unsure{envers} et contre
tout?)", de faire correspondre à $\tau_{X_0}$ l'opération
$\tau_I$\inserted{\unsure{donné|donnée} par (27)}, qui sur
chaque $T_j$ serait \unsure{égal|égale} à $\tau'_j$
(et sur $V_0$, bien sûr, coïncide avec $\tau_{X_0}$),
$\tau'_j(z,t)=(z,-t)$,
\deleted{donc} l'ens des pts fixes de $\tau_I^{\unsure{2}}$
étant formé des $g$ cercles \unsure{médians} $S_j\times\{0\}$
des $g$ tubes $T_j$. On aura par (41)
$\tau_I=\rho_I\tau_g$, donc $\tau_I$ commute
\unsure{également} à $\mathbb D_n$, puisque $\tau$ et
$\rho_I=\prod_{i\in I}\rho_i$ y commutent.
Mais il \unsure{ne|me} semble pas \unsure{important} pour le moment
quelle convention nous \unsure{adoptons}]

On peut donc dire aussi que le groupe $\mathbb D_{2g}^\sim$
opère sur $\mathcal H$ --- cette opération prolongeant celle
de $\tau_g$\deleted{\unsure{par}}, identifié maintenant à un
élément de $\mathbb D_{2g}^\sim$, i.e. à l'antipodisme dans
$\mathbb D_{2g}^\sim$ --- et le produit $\frac12$ direct
\[
\tag{51}
\mathbb D_{2g}^\sim.\mathcal H\supset\mathcal G
=\deleted{\unsure{\mathcal H}}\{1,\tau_g\}.\mathcal H
\]

\margin{465}

\scanpage{467}
\mspage{236}

opère sur $X_g$.

Il faudrait maintenant, dans cette voie :

1) Expliciter l'action extérieure de ce produit $\frac12$ direct sur le groupe fondamental $\unsure{\pi_g|\Pi_g}$, \unsure{avec} une attention toute particulière à l'action du groupe $(\mathbb Z/2\mathbb Z)^3\subset\mathbb D_{2g}^{\sim}$ qui stabilise un des cylindres $T_j$ --- qui dans l'espace euclidien de dim 3 s'interprète comme le groupe des \unsure{chgts} de signe relatif \unsure{au|un} syst. d'axes orthonormé correspondant.

2) Étudier l'action \inserted{sur $\unsure{\pi_g|\Pi_g}$} du groupe de Teichmüller (plus ou moins) ``spécial'' de chacun des $T_j$, ou l'action d'un groupe analogue d'un ensemble plus grand obtenu en lui rajoutant une ``lanière'' \inserted{$L$} d'où un tore à 1 trou, et son groupe de Teichmüller spécial, qui \unsure{s'introduisent} de façon naturelle$^*$) \unsure{voire|voir} l'ensemble encore plus grand obtenu en mettant également la lanière antipodique $L'$ (cet ensemble se présente aussi comme un cylindre $(L\cup L')$, ou ``buse'', où on aurait mis un tube $(T_j)$ en travers, et a la structure topologique d'un tore à deux trous.

\figure{[Grand dessin à gauche du deuxième alinéa : contour ovale presque circulaire; deux points noirs sur ce contour, l'un en haut, l'autre en bas. Deux traits courbes pleins descendent du point supérieur vers une petite ouverture circulaire à gauche; la bande ainsi formée est marquée $L$. Deux traits pleins repartent de cette ouverture vers le point inférieur; la bande inférieure est marquée $L'$. Du côté droit, des traits courbes pointillés relient les deux points du contour à une seconde petite ouverture circulaire. Deux traits pointillés obliques joignent les ouvertures de gauche et de droite et encadrent l'indication $T_j$. À droite de ce tube se trouvent une petite marque noire et un petit ovale hachuré. Aucun sens de flèche n'est indiqué.]}

\margin{$^*$ C'est essentiellement un $\mathrm{SL}(2,\mathbb Z)$ --- plutôt une extension centrale \unsure{remarquable} de $\mathrm{SL}(2,\mathbb Z)$ par $\mathbb Z$, qui \unsure{rappelle} celle de $\mathrm{SL}(2,\mathbb R)$ par $\mathbb Z$\ldots\ (rev. universel de $\mathrm{SL}(2,\mathbb R)$)}

\figure{[Petit dessin sous le précédent : deux grands contours ovales verticaux en partie superposés, réunis en haut et en bas; une partie du contour intérieur est pointillée. Un tube transversal, marqué $T_j$, joint deux petites ouvertures circulaires à gauche et à droite. Les bords visibles du tube sont pleins; des segments cachés près des ouvertures sont pointillés. Aucun sens de flèche n'est indiqué.]}

\margin{466}

\scanpage{468}

Il est possible qu'il faudra considérer \unsure{des|les} \unsure{pris|pas} a deux, pour étudier les relations entre \inserted{les} éléments de $\mathcal T_g$ provenant des ensembles précédents\ldots\ Un travail amusant sera de se débrouiller pour écrire les générateurs du groupe $\unsure{\pi_g|\Pi_g}$, et surtout la fameuse relation, dans une disposition géométrique relative des ``anses'', qui est une disposition ``\unsure{panachée}'' --- et non plus \uline{sagement} à la queue-leu-leu !

\deleted{Je vois \unsure{apparaître}}

En attendant d'entrer ainsi dans le vif de la structure du groupe de Teichmüller, je vais déjà essayer de décrire des générateurs et relations pour le sous-groupe intéressant qu'on vient d'écrire, $\mathbb D_{2g}^{\sim}\cdot H$.
Je vais prendre les sempiternels générateurs $\varepsilon_0,\varepsilon_1$ de $\mathbb D_{2g}$
\[
\varepsilon_0^2=\varepsilon_1^2=1,\qquad (\varepsilon_0\varepsilon_1)^{2g}=1
\]
et y joindre $\tau=\tau_g$, et \deleted{\illegible{1}} $\rho_0\tau=\tau'$
($\rho_0=\rho_{j_0}$, $j_0$ pt marqué de $I=J/\tau$)
satisfaisant
\[
\tau^2=\tau'^2=1,\qquad
\deleted{[\tau,\varepsilon_0]=\unsure{[\varepsilon_1,\tau]|[\tau,\varepsilon_1]}=}
\]
\[
(\tau\varepsilon_0)^2=(\tau\varepsilon_1)^2=1\quad
\deleted{\unsure{(\tau'\varepsilon_0)^2=(\tau'\varepsilon_1)^2}=1}
\]
\[
(\varepsilon_1\tau')^2=1
\qquad \varepsilon_1\ \text{commute à}\ \tau'
\]

\margin{\uline{NB} $\varepsilon_0$ bouge \unsure{l'élément de repère}, $\varepsilon_1$ \unsure{le fixe} et \illegible{1} \unsure{en somme}, donc $\varepsilon_1(j_0)=j_0$}

\margin{\unsure{expriment} la \unsure{commutativité|commutation} de $\tau$ \deleted{\illegible{1}} à $\varepsilon_0,\varepsilon_1$}

\margin{467}

\scanpage{469}
\mspage{237}

\[
((\varepsilon_0\varepsilon_1)^g\tau')^2=1
\]
\margin{(l'involution $(\varepsilon_0\varepsilon_1)^g$ commute à $\tau'$)}

sauf erreur, ça fait un ensemble de générateurs et relations pour $\mathbb D_{2g}^{\sim}\cdot\mathcal H$ --- en résumé $\varepsilon_0,\varepsilon_1,\tau,\tau'$
\[
\text{(52)}\quad
\left\{
\begin{array}{l}
\varepsilon_0^2=\varepsilon_1^2=\tau^2=\tau'^2=1\\
(\tau\varepsilon_0)^2=(\tau\varepsilon_1)^2=(\tau'\varepsilon_1)^2=1\\
(\varepsilon_0\varepsilon_1)^{2g}=1\\
(\tau'(\varepsilon_0\varepsilon_1)^g)^2=1
\end{array}
\right.
\]

C'est peut-être pas très astucieux comme choix de générateurs, en ce sens que $\varepsilon_0,\varepsilon_1$, ne sont pas du tt sur le même pied que $\tau,\tau'$ --- ce ne sont pas des \uline{anti}-\unsure{involutions}. Il vaudrait mieux prendre $\varepsilon'_0=\tau\varepsilon_0$, $\varepsilon'_1=\tau\unsure{\varepsilon_1}$ de façon à obtenir
\[
\text{(52bis)}\quad
\left\{
\begin{array}{l}
\varepsilon_0^{\prime 2}=\varepsilon_1^{\prime 2}=\tau^2=\tau'^2=1\\
(\tau\varepsilon'_0)^2=(\tau\varepsilon'_1)^2=(\tau'(\tau\varepsilon'_1))^2=1\\
(\varepsilon'_0\varepsilon'_1)^{2g}=1\\
(\tau'(\varepsilon'_0\varepsilon'_1)^g)^2=1
\end{array}
\right.
\]
ce sont essentiellement ``les mêmes'' relations sauf la dernière de la deuxième ligne.

Peut être ce petit jeu avec le tout petit groupe $\mathbb D_n^{\sim}$ opérant sur $\mathcal H$ est un peu une \unsure{amusette} --- le groupe $\mathcal H\simeq\mathbb Z^g$ dans $\mathcal T_g$ suggère beaucoup la situation

\margin{468}

\scanpage{470}

d'un tore maximal dans un groupe $\frac12$ simple, on a vraiment envie d'en voir une caractérisation intrinsèque dans $\mathcal T_g$ --- ou dans $\unsure{\mathcal G_g|\mathcal C_g}$, ce qui est possible puisque \deleted{\illegible{1}} en tant que groupe de transformations topologiques \unsure{effectives|effectif}, il laisse fixe les points de $V_0\subset X_g$ --- p.ex. les deux pôles, qu'on peut prendre comme pts \unsure{bases|base} pour construire rev. universels et groupes \unsure{fondamental|fondamentale}. Ce sont peut-être les sous-groupes \unsure{abéliens|abélien}-libres maximaux.

On aimerait étudier le normalisateur --- est-ce exactement ce qui provient des autom de $X_g-\bigcup\mathring{T}_i\simeq V_0$ ?

les $\mathrm{SL}(2,\mathbb Z_{\illegible{1}})$ associés aux ``lanières'' à travers les \unsure{tubes|trous} \inserted{$T_j$} jouent-ils un rôle analogue à celui des groupes $\mathrm{SL}(2,K)$ ou $\unsure{\mathrm{GP}(1,k)|\mathrm{GL}(1,k)}$ ``de rang 1'' dans la théorie des groupes algébriques réductifs ? Si Jean \illegible{1} ``\unsure{mord}''

\margin{469}
